
Language models deployed in applications requiring reliable uncertainty estimates must produce confidence scores that discriminate between correct and incorrect predictions. Prior work has documented that RLHF-tuned models exhibit systematic overconfidence compared to base models \citep{Tian2023Calibration}, and post-hoc temperature scaling has been shown to reduce Expected Calibration Error (ECE) in neural networks \citep{Guo2017Calibration}. Recent methods such as DACA \citep{Luo2025DACA} and CCPS \citep{Khanmohammadi2025CCPS} report ECE improvements of 15-55\% on large language models.

However, ECE measures whether predicted probabilities match observed frequencies, not whether confidence scores effectively rank predictions by correctness probability. A model can achieve low ECE through probability rescaling while having degraded ability to distinguish correct from incorrect predictions. This distinction motivates examining whether RLHF affects the discriminative quality of confidence signals, measured by AUROC rather than calibration metrics alone.

This work tests the hypothesis that RLHF instruction tuning degrades discriminative confidence quality through a mechanism distinct from scalar miscalibration. Specifically, we hypothesize that preference optimization inflates logit margins including for incorrect predictions, fundamentally reshaping the confidence-correctness relationship rather than merely shifting its scale.

We make four contributions:

\begin{enumerate}
\item We measure discriminative degradation using AUROC for margin-based correctness prediction across two model families (Qwen and Mistral), finding consistent decreases of 2-4 percentage points in instruction-tuned models.
\end{enumerate}

\begin{enumerate}
\item We identify the mechanism as disproportionate margin inflation for incorrect predictions, with inflation ratios of 3-17x depending on model family.
\end{enumerate}

\begin{enumerate}
\item We use percentile-normalized logistic regression to separate scale effects from shape effects, finding that slope attenuation persists after normalization.
\end{enumerate}

\begin{enumerate}
\item We apply Murphy's Brier score decomposition to characterize the distortion type, finding refinement degradation consistent with geometric rather than scalar distortion.
\end{enumerate}
